\begin{frame}{From Concepts to Application Design}

    \begin{center}

        \begin{tikzpicture}[
            font=\small,
            concept/.style={
                draw=WorkshopBlue,
                fill=WorkshopLightBlue,
                rounded corners=4pt,
                minimum width=2.8cm,
                minimum height=0.75cm,
                align=center
            }
        ]

            % Concepts
            \node[concept] (functions) {Functions};
            \node[concept, right=0.35cm of functions] (modules) {Modules};
            \node[concept, right=0.35cm of modules] (oop) {OOP};
            \node[concept, right=0.35cm of oop] (storage) {JSON / Storage};

            % Arrow
            \node[
                below=0.55cm of modules,
                font=\large\bfseries
            ] (arrow) {
                $\Downarrow$
            };

            % Application
            \node[
                draw=WorkshopBlue,
                fill=WorkshopGreen,
                rounded corners=6pt,
                minimum width=7.5cm,
                minimum height=1.1cm,
                below=0.45cm of arrow,
                font=\Large\bfseries
            ] (application) {
                Student Management System
            };

        \end{tikzpicture}

    \end{center}

    \vspace{0.35cm}

    \begin{itemize}

        \item Functions organize reusable logic

        \item Modules organize the code

        \item OOP models the student and related entities

        \item JSON provides persistence

    \end{itemize}

    \vfill

    \begin{center}
        \large
        \textbf{Individual concepts become building blocks of an application.}
    \end{center}

\end{frame}